
\subsection*{7.13. Relations avec la cohomologie \(\ell\)-adique}
Rappelons (SGA 5 VII) que pour tout Module localement libre de type fini \(\underline{E}\) sur un schéma \(T\), et tout nombre premier distinct des caractéristiques résiduelles de \(T\), on définit des classes de Chern \(\ell\)-adiques (pour \(i\geq 0\))
\[
  c^i(E)\in H^{2i}(T,\mathbb Z_\ell(i)),
\]
satisfaisant à une formule d'additivité
\[
  c(E)=c(E')c(E'')
\]
lorsque \(E\) est extension d'un module localement libre \(E'\) par un module localement libre \(E''\). De cette additivité on conclut que la classe de Chern totale \(c(E)=\sum_i c^i(E)\) provient d'un homomorphisme
\begin{equation}\tag{7.13.0}
  c:K^\bullet(T)\longrightarrow
  \widehat G\bigl(H^{2*}(T,\mathbb Z_\ell(*))\bigr)
  =
  1+\prod_{i\geq 1}H^{2i}(T,\mathbb Z_\ell(i)),
\end{equation}
(où l'objet but est un groupe par le produit, comme expliqué dans V 1, et où le \(K^\bullet(T)\) est le \(K^\bullet\) «~\emph{naïf}~»). On en déduit des applications canoniques
\begin{equation}\tag{7.13.1}
  c^i:K^\bullet(T)\longrightarrow H^{2i}(T,\mathbb Z_\ell(i)).
\end{equation}

Ceci posé, soit \(X\) un schéma \emph{propre} et plat sur \(S\), à \emph{fibres régulières}, et supposons \(S\) \emph{hensélien}. On suppose \(\ell\neq\) car. résiduelle de \(S\). On peut alors former le diagramme
\begin{equation}\tag{7.13.2}
\begin{tikzcd}[column sep=large,row sep=large]
K^\bullet(X_t) \arrow[r,"c^i"] \arrow[d,"\sigma"'] &
H^{2i}(X_t,\mathbb Z_\ell(i))\\
K^\bullet(X_s) \arrow[r,"c^i"] & H^{2i}(X_s,\mathbb Z_\ell(i)) \arrow[u,"\sigma^{-1}"']
\end{tikzcd}
\;,
\end{equation}
où les flèches horizontales sont les applications «~classe de Chern~» relatives à \(X_t\) et \(X_s\) respectivement, et où \(\sigma^{-1}\) est l'homomorphisme composé
\[
  H(X_s)\xleftarrow{\substack{i^*\\\sim}}H(X)\xrightarrow{j^*}H(X_t);
\]
